\documentclass{article}
\usepackage{alexconfig}
\title{Basic Algebraic Geometry Section 1 Part 1 Notes}

\begin{document}
\maketitle
\

\section{Plane Curves}

\begin{definition}[Fields]
	\

	A \textbf{field} is a ring where every element is a unit (has a multiplicative inverse) besides $0$
\end{definition}

\begin{definition}[Field Extension]
	\

	Given a field $k$, another field $F$ is called a \textbf{extension field} of $k$ if the operations in $F$ restricted to $k$ are the same as the operations of $k$. So, for $a,b,c \in k$, if $$ab = c$$in $F$, then $$ab = c$$in $k$.

	The pair of fields $F, k$ is called a \textbf{field extension} and denoted $$\frac{F}{k}$$

	In addition, we can think of $F$ as a vector space over $k$, and the basis of $F$ as a $k$-vector space is the \textbf{degree of the field extension} and is denoted $[F:k]$.
\end{definition}

\begin{definition}[Affine Space]
	\

	Given a field $k$, the \textbf{affine space} denoted $\mathbb{A}^n$ is the cartesian product $k^n$. So, for $a_1 \dots a_n \in k$, elements of $\mathbb{A}^n$ look like $$(a_1, \dots, a_n)$$
\end{definition}

\begin{definition}[Algebraic Plane Curve]
	\

	An \textbf{algebraic plane curve} is the set of points which satisfy the equation $$f(x,y) = 0$$where $f$ is a nonconstant polynomial, $x, y \in k$ for some field $k$, and the coefficients of $f$ are also in $k$.

	Then, $f \in k[x,y]$.

	A curve of degree $2$ is called \textbf{conic}, a curve of degree $3$ is called \textbf{cubic}
\end{definition}


\begin{definition}[Unique Factorization Domain]
	\

	An integral domain $R$ (ring with no zero divisors) is a \textbf{unique factorization domain} if for all $r \in R$, $r$ can be written as a finite product of irreducible elements, $$r = r_1 r_2  \dots r_n$$where an irreducible element $r_1 \in R$ is an element with no multiplicative inverse.

	Uniqueness is up to multiplication by units, because for some factor $r_1$, if we have unit $u$, we can just write $u u ^{-1} r_1$, and this should not be counted as "unique".
\end{definition}

\begin{example}
	\

	The polynomial ring $k[x,y]$ is a UFD.
\end{example}

\begin{proposition}
	\

	If some $f \in k[x,y]$ has factors $$f = f_1 f_2 \dots f_n$$then the zeros of $f$ are the union of all the zeros of the individual factors $f_i$.
\end{proposition}

\begin{customproof}
	\

	If $(x_i, y_i)$ is a zero for some $f_i$, then $f(x_i,y_i) = 0$, so $$f(x_i, y_i) = f_1 \dots f_{i-1} \cdot 0 \cdot f_{i+1} \dots f_n = 0$$so $(x_i, y_i)$ is a zero for $f$ as well.
\end{customproof}

\begin{lemma}
	\

	Let $k$ be an arbitrary field, $f \in k[x,y]$ be an irreducible polynomial, $g \in k[x,y]$ be some other polynomial. If $g$ is not divisible by $f$ then the system of equations $f(x,y) = g(x,y) = 0$ has only a finite number of solutions.
\end{lemma}

\begin{customproof}
	\

	Consider the ideal $(f)$, and take $k[x,y]/(f)$. In this new ring, zeros of elements are when the function agrees with $f$. Then, since we know $g$ is not divisible by $f$, it cannot be $(f)$, so it cannot be identically zero, so it must only agree on finitely many points.
\end{customproof}

\begin{definition}[Algebraically Closed]
	\

	A field $k$ is algebraically closed if every nonconstant polynomial with one variable with coefficients in $k$ has a root in $k$. So, for all nonconstant $f \in k[x]$, $f(a) = 0$ for some $a \in k$.
\end{definition}

\begin{example}
	\

	$\mathbb{R}$ is not algebraically closed because $x^2 + 1 = 0$ has no valid solutions.
\end{example}

\begin{lemma}
	\
	A field $k$ being algebraically closed is equivalent to the following statements:

	\begin{enumerate}
		\item The only irreducible polynomials are those of degree one
		\item Every polynomial is a product of first degree polynomials
		\item Polynomials of prime degree have roots
		\item The field has no proper algebraic extension
		\item The field has no proper finite extension
		\item For any two polynomials $f, g \in k[x]$, if $f,g$ are coprime, they do not share a common root.
	\end{enumerate}
\end{lemma}

\section{Rational Curves}
\

\begin{example}
\

The curve $$y^2 = x^2 + x^3$$can be expressed via two rational functions of one parameter.
\end{example}

\begin{customproof}
\

Observe that any line of the form $y = tx$ has exactly one point of intersection with our curve, besides the origin. To show this, let's plug in: $$(tx)^2 = x^2 + x^3$$so$$x^2(t^2 - x - 1) = 0$$Our $x^2$ root corresponds to the intersection point at the origin, $(0,0)$. Then, we have $$t^2 - x - 1 = 0$$so$$x = t^2 - 1$$Then substituting back into our equation, we get that $$y = t(t^2-1)$$

There is actually a $1$-$1$ association between points on the curve (besides $(0,0)$) and 
\end{customproof}




\end{document}
